\chapter{What transfers, what is new, what is noise}
\label{ch:gap-map}

\section*{What this chapter is for}

You are not starting from zero and you are not starting from where you think.
Most of what a senior backend or platform engineer already knows transfers
directly, under different names. A smaller set of things is genuinely new. A
third set looks essential, is studied first by almost everybody making this
move, and is very unlikely to be asked about at all.

This chapter is the map of those three sets. It is the first thing to read and
the thing most likely to save you a month.

\section{What transfers, and what it is called now}

The useful form of this is a mapping rather than a list, because the point is
not that your skills are \emph{adjacent} to the new ones. In several cases they
are the same skill with a different noun.

\begin{kittable}{@{}XX@{}}
\textbf{What you already do} & \textbf{What it is called here} \\
\midrule
Designing an API for a caller you do not control & Designing a tool schema for a model you do not control \\
Integration and regression tests & Evals \\
Rate limits, retries, circuit breakers, fallbacks & Gateway resilience across model providers \\
Logs, metrics, traces & Token, cost and latency tracing \\
Input validation at a trust boundary & Treating model output, and retrieved documents, as untrusted \\
Capacity planning & Context budgeting \\
Schema migration under live traffic & Prompt and model version migration under live traffic \\
\end{kittable}

The right-hand column is not a translation exercise. The engineering underneath
is the engineering you have: \reported{Python or TypeScript, API design, testing
and continuous integration, deployment and monitoring all transfer directly, and
the engineering foundation is the asset.}{field-guide-backend}

\judgement{The row worth dwelling on is the second. If you have ever argued
that a test suite which has never failed is not evidence, you already hold the
central idea of evaluation, and you hold it more firmly than someone who
arrived at it through a library's documentation.}

\section{What is genuinely new}

Four things, and only the last is hard.

\textbf{Calling a model well.} Structured output, function calling, the
difference between a system instruction and a user turn, what a context window
costs. A week, unhurried.

\textbf{Prompting as an engineering discipline.} Not writing clever prompts:
versioning them, testing them, and treating a change to one as a change to
behaviour that requires the same review a code change does.

\textbf{Retrieval.} What it is for, how it fails, and when it is the wrong
answer.

\textbf{Non-determinism, and what it does to everything you know about
testing.} The clearest statement of the gap is a question rather than a topic,
and it was put by a backend engineer making this exact move: \emph{how do you
even know something is broken when the output is different every
time?}\source{field-guide-backend}

That question is the whole of the fourth item. Every familiar instrument
assumes repeatability. A failing assertion tells you something changed. A green
suite tells you nothing changed. Neither statement survives here, and the
machinery that replaces them \dash{} baselines, absence assertions, anchored
rubrics, statistical gates \dash{} is the subject of Chapter~\ref{ch:evals}.

\reported{Evaluation is reported as the steepest part of the curve for
engineers arriving from backend work, and as the thing they do not
have.}{field-guide-backend} \judgement{It is also the thing that, once you have
it, distinguishes you most sharply from a candidate who arrived from data
science, because they tend to have the statistics and not the instinct that a
check nobody has watched fail is not known to be checking anything.}

\section{What is noise}

This list is the reason to read this chapter before starting to study.

\begin{kittable}{@{}lXl@{}}
\textbf{Topic} & \textbf{Why it looks essential} & \textbf{Reality} \\
\midrule
Fine-tuning, LoRA, PEFT & every course leads with it & niche\source{field-guide} \\
Transformer internals & it is what \emph{AI} means to most people & research-adjacent only\source{field-guide} \\
Attention arithmetic, KV cache & it is teachable and feels rigorous & rarely asked\source{field-guide} \\
Mixture-of-experts, quantisation & it is in the news & rarely asked\source{field-guide} \\
\end{kittable}

The corpus states the rule plainly: \reported{if the posting does not mention
fine-tuning, you are unlikely to be asked about it.}{field-guide} The same holds
for the rest of the list.

\begin{kittrap}
The trap is not that this material is worthless. It is that it is the only part
of the subject with a settled curriculum, so it is the part that has courses,
and a course is what an engineer reaches for when starting something new. You
can spend a month becoming able to derive scaled dot-product attention and walk
into a room where the first question is what happens to your retrieval quality
when a document is longer than the chunk size.
\end{kittrap}

There is a second kind of noise: material that is genuinely useful but has been
renamed, so it looks like a new subject to learn. Several currently-promoted
topics turn out, on inspection, to be retrieval, prompt design and evaluation
under a newer word. Chapter~\ref{ch:noise} names the ones that were traceable
when this was written, and says how to test a new one yourself.

\section{Where to put the time}

\begin{kitmethod}
In order, and stop when you run out of weeks rather than reordering:

\textbf{1.} Build one small thing end to end that calls a model, uses a tool,
and is deployed somewhere with a URL. A day or two. It converts every abstract
question into a concrete one.

\textbf{2.} Put an evaluation harness on it \dash{} however small. This is the
highest-yield week available to you, for reasons Chapter~\ref{ch:portfolio}
gives.

\textbf{3.} Learn retrieval properly, including its failure modes, because it
is the single most-asked surface.

\textbf{4.} Do the cost and latency arithmetic on the thing you built, out
loud, on paper.

\textbf{5.} Only then read about anything in the noise table, and only if a
posting you are actually pursuing mentions it.
\end{kitmethod}

\begin{drill}
Write down, in one page, the architecture of a system you have shipped. Now
annotate it with the right-hand column of the transfer table: mark every place
where you already solved the problem that the new vocabulary is about. Keep the
page. It is the raw material for the project deep-dive in
Chapter~\ref{ch:deep-dive}, and it is the fastest way to stop feeling like a
beginner in a field where you are not one.
\end{drill}

\section*{What this chapter does not claim}

It does not claim the noise table is permanent. A posting that asks for
fine-tuning means the role does fine-tuning, and then it is not noise for that
role \dash{} the table is about the default, not about every case. Re-check it
against the postings in front of you rather than against this page.

It does not claim that two to three months is enough to make the transition.
That figure circulates and traces to a single source; it is refused in
Appendix~\ref{app:refused} along with the rest of that family. What this chapter
claims is narrower: that the order above is better than the order almost
everybody uses, which is to start with the mathematics.
